% Préambule du rapport de stage — packages et réglages généraux.
\usepackage[french]{babel}
\usepackage[T1]{fontenc}
\usepackage{lmodern}
\usepackage[a4paper, margin=2.5cm]{geometry}
\usepackage{graphicx}
\usepackage{xcolor}
\usepackage{hyperref}
\usepackage{fancyhdr}
\usepackage{titlesec}
\usepackage{setspace}
\usepackage{enumitem}
\usepackage{longtable}
\usepackage{array}
\usepackage{csquotes}

% Couleurs institutionnelles simples, réutilisées dans les titres et la page de garde.
\definecolor{bleuinstitutionnel}{RGB}{15, 40, 90}
\definecolor{grisdiscret}{RGB}{90, 90, 90}

\onehalfspacing

% Style des titres de chapitre et de section.
\titleformat{\chapter}[display]
  {\normalfont\huge\bfseries\color{bleuinstitutionnel}}
  {\chaptertitlename\ \thechapter}{20pt}{\Huge}
\titleformat{\section}
  {\normalfont\Large\bfseries\color{bleuinstitutionnel}}{\thesection}{1em}{}
\titleformat{\subsection}
  {\normalfont\large\bfseries}{\thesubsection}{1em}{}

% En-têtes et pieds de page du corps du document.
\pagestyle{fancy}
\setlength{\headheight}{15pt}
\fancyhf{}
\renewcommand{\headrulewidth}{0.4pt}
\fancyhead[L]{\small\textit{\leftmark}}
\fancyhead[R]{\small Rapport de stage}
\fancyfoot[C]{\thepage}

\hypersetup{
  colorlinks=true,
  linkcolor=bleuinstitutionnel,
  citecolor=bleuinstitutionnel,
  urlcolor=bleuinstitutionnel,
  pdftitle={Rapport de stage},
  pdfauthor={Diogo Balde},
}
